% ============================================================
% Chapter 4: Proposed Method: Specification and Design
%   Section: Enrichment
% Included from chapters/04-proposed-method/proposed-method.tex
% Source: docs/arquitectura-tecnica.md section 2.3; docs/decisions/
% inference.md; src/processors/nlp/.
% ============================================================

\section{Enrichment}
\label{sec:design-enrichment}

Every article that comes out of the extraction cascade
(Subsection~\ref{subsec:scraping-cascade}) is turned from plain text
into an enriched record by a stack of language-processing steps. Each
runs in the article's language, detected from stopword frequencies,
with English and Spanish word lists for every heuristic.

\begin{itemize}
    \item \textbf{Keywords} by YAKE, an unsupervised statistical method
    that needs no training data \cite{campos2020yake}, with the
    article's language's stopwords.
    \item \textbf{Entities} by GLiNER, a named-entity recogniser that
    takes the entity types as input rather than fixing them at training
    time \cite{zaratiana2024gliner}: people, organisations, places,
    countries and others. The model reads at most 384 tokens, so a long
    article is read in sentence-bounded chunks; until 16~October it was
    silently truncated to its opening (Section~\ref{sec:dev-inference}).
    \item \textbf{Candidate claims}: each sentence is scored by the
    presence of structural signals of a checkable fact --- named
    entities, figures, dates, units, reporting verbs such as
    ``announced'' or ``confirmed'', quotations, a minimum length --- and
    kept above a threshold. The score measures whether a sentence has
    the \emph{form} of a verifiable fact, not whether it is important;
    a sentence must also end as a sentence, which removed headlines,
    captions and fragments cut at an abbreviation.
    \item \textbf{Topics}: the article's embedding is compared with
    descriptions of 23 topics, organised in five groups. Each topic gets
    a confidence, and a topic counts when its confidence clears a
    threshold (0.35). A softmax over the same similarities is also
    computed, but over 23 near-equal cosines it is close to uniform and
    carries no signal; nothing downstream uses it.
    \item \textbf{Sentiment} by a multilingual XLM-RoBERTa classifier
    \cite{barbieri2022xlmt}, in full precision: an int8 version was
    measured on 336 real texts and rejected because it changed the
    label in 8--14\% of them, mostly from neutral to positive, which
    would have quietly loosened the admission gate.
    \item \textbf{Editorial quality}: readability, objectivity,
    constructiveness, hopefulness, social impact and inspirational
    value, from curated word lists per language and standard
    readability formulas.
    \item \textbf{An embedding} of the whole article by BGE-M3, a
    multilingual embedding model \cite{chen2024bgem3}. The same
    1,024-dimensional space underlies duplicate detection, topic
    classification, the mission screen, claim deduplication and
    evidence ranking.
\end{itemize}

The three transformer models run in the inference service
(Section~\ref{sec:architecture}). Each step fails differently by
design: if entity extraction fails, the claim scorer loses one weak
signal and carries on; if sentiment or embeddings fail, the article
fails, because a silently neutral sentiment or a zero vector would
corrupt an admission decision without anyone noticing.
